\begin{proposition}[Meridian quotient and rational incidence trees]\hypertarget{tag-02YG}{}\label{night:P4R.4}\label{stmt:P4R.4}
\StatementRecord{night:P4R.4}{proved}{02YG}{\SeriesRef{night:P4.13}, \SeriesRef{night:P4R.1}, \SeriesRef{night:P4R.3}}

In the setting of \SeriesRef{night:P4R.1}, choose a positive small-loop
monodromy $A$ on $\Lambda=H_1(T^4;\mathbb Z)$.
Let $t$ lift the positive base generator, with the convention
$tvt^{-1}=A(v)$ for $v\in\Lambda$.
For a suitable $\lambda\in\Lambda$, a positive
meridian of the reduced support is represented by
$\mu=\lambda t^m$ in $\Lambda\rtimes_A\langle t\rangle$.
The natural small-tube comparison gives
\[
 \pi_1(S)\simeq
 (\Lambda\rtimes_A\langle t\rangle)/\langle\!\langle\mu\rangle\!\rangle,
 \qquad
 \pi_1(S)/\operatorname{im}\Lambda\simeq\mathbb Z/m,
 \qquad
 b_1(S)=\dim_{\mathbb Q}\ker(A-I).
\]
For the support in \SeriesRef{night:P4R.3}, a singular-locus incidence
graph that is a tree of rational normalization
components gives $e(S)\ge3$.
\end{proposition}
\begin{proof}

\begin{enumerate}[label=\arabic*.,leftmargin=*]
\item Use the nested neighbourhoods in the proof of
\SeriesRef{night:P4R.1}. Choose an open-star
neighbourhood $V_0$ retracting onto $S$, a small tube
$U_0\subset V_0$, a second open-star neighbourhood
$V_1$ with $S\subset V_1\subset U_0$, and a smaller
tube $U_1\subset V_1$. Properness supplies both tube
inclusions. The punctured-disc fibrations give an
isomorphism
\[
 \pi_1(U_1\setminus S)\longrightarrow
 \pi_1(U_0\setminus S):
\]
the fibres can be identified by radial parallel
transport and both base fundamental groups are
generated by the same positive loop. As in the proof of
\SeriesRef{night:P4R.1},
the kernels of the maps to $\pi_1(U_i)$ are the
normal closures of the positive smooth-support
meridian. These meridians correspond under the
displayed isomorphism. Hence
$\pi_1(U_1)\to\pi_1(U_0)$ is an isomorphism.

\item This isomorphism factors through $\pi_1(V_1)$,
which is identified with $\pi_1(S)$ by its
retraction. Thus $\pi_1(S)\to\pi_1(U_0)$ is
surjective. The restriction to $U_0$ of the
retraction of $V_0$ proves injectivity. The inclusion
$S\hookrightarrow U_0$ and the restricted retraction are
inverse on fundamental groups. Together with the meridian
kernel calculation, this gives the asserted quotient
presentation of $\pi_1(S)$.

\item Projection to the cyclic factor sends $\mu$ to
$m$. After dividing by the image of $\Lambda$,
the remaining presentation is therefore exactly
$\langle t\mid t^m=1\rangle$. In the rational
abelianization the presentation becomes
\[
 \frac{(\Lambda/(A-I)\Lambda)\otimes\mathbb Q
       \ \oplus\ \mathbb Qt}
      {\mathbb Q(\overline\lambda+mt)}.
\]
The single displayed relation has nonzero
$t$ coefficient and reduces the dimension by
one. Explicitly, with
$V=(\Lambda/(A-I)\Lambda)\otimes\mathbb Q$, the map
$(v,c)\mapsto v-c\overline\lambda/m$ induces an isomorphism
of this quotient with $V$, whose inverse sends $v$ to
$[(v,0)]$. This proves the asserted formula for $b_1$.

\item In the rational-tree case, \SeriesRef{night:P4R.3}
gives $\pi_1(S)\simeq\pi_1(C)$ for $C=\operatorname{Sing}(S)_{\rm red}$.
Write $\nu_C:\coprod_{j=1}^c\mathbb P^1_j\to C$ for its
normalization. Replace each identification of a finite fibre
of $\nu_C$ by a star joining its preimages. Collapsing these
contractible stars gives $C$ up to homotopy; collapsing the
sphere components induces an isomorphism on fundamental
groups with the incidence graph. Van Kampen therefore gives
$\pi_1(C)=1$ and $\pi_1(S)=1$.
The cyclic quotient then forces $m=1$, and the
$b_1$ formula gives $\det(I-A)\ne0$.
\SeriesRef{night:P4.13} gives
$e(S_{\rm reg})=\det(I-A)\ge0$; this nonzero integer
is at least one.
Let $P=\{p\in C:\#\nu_C^{-1}(p)>1\}$ and put
$r_p=\#\nu_C^{-1}(p)$ for $p\in P$. The tree has
$c+\#P$ vertices and $\sum_{p\in P}r_p$ edges, so
$\sum_p(r_p-1)=c-1$. Thus
\[
 e(C)=2c-\sum_p(r_p-1)=c+1\ge2.
\]
Unibranch singular points have one preimage and contribute
zero to the subtraction. Finally
$e(S)=e(S_{\rm reg})+e(C)\ge3$.

\end{enumerate}
\end{proof}